\section{Fields}

Throughout the semester, we will employ the following conventions:
\begin{itemize}
    \item Unless otherwise specified, all rings will be assumed to be unital and commutative.
    \item A \emph{ring homomorphism} is a map of rings $\varphi : R \to S$, such that
    \begin{enumerate}
        \item $\varphi(x+y) = \varphi(x) + \varphi(y)$, for all $x, y \in R$.
        \item $\varphi(x \cdot y) = \varphi(x) \cdot \varphi(y)$, for all $x, y \in R$.
        \item $\varphi(1_R) = 1_S$
    \end{enumerate}
\end{itemize}

\begin{definition}
    Let $R$ be a ring. An element $x \in R$ is called a \emph{unit} or \emph{invertible}, if there exists an element $y \in R$ such that $x \cdot y = 1$.
\end{definition}

The set of all units of $R$ form a group under multiplication and will be denoted $R^*$.

\begin{example}{\ }
    \begin{enumerate}
        \item $\Z^* = \{ -1, 1 \}$
        \item $\Q^* = \Q \setminus \{0\}$
        \item For $n \in \N$, it holds that $(\Z_n)^* = \{ m \in \Z_n : \gcd(m,n) = 1 \}$.
    \end{enumerate}
\end{example}

\begin{definition}
    A \emph{field} is a ring $\K$ such that $\K^* = \K \setminus \{0\}$.
\end{definition}

\begin{example}{\ }
    \begin{enumerate}
        \item $\Q \subsetneq \R \subsetneq \C$
        \item $\Z_p$, where $p \in \P$.
    \end{enumerate}
\end{example}

\begin{lemma}
    A ring $R$ is a field if and only if the only ideals of $R$ are $\{0\}$ and $R$.
\end{lemma}

\begin{proof}
    Let $I \subseteq R$ be an ideal. If $I$ contains a unit then $I = R$.

    ``$\implies$'': If $I \subseteq R$ is a non-zero, non-trivial ideal, then $I$ contains a unit, which implies $I = R$.

    ``$\impliedby$'': If $R$ is not a field, then there exists some non-zero, non-unit element $x \in R$. But then the ideal $Rx$ satisfies
    \[
        \{ 0 \} \subsetneq Rx \subsetneq R. \qedhere
    \]
\end{proof}

\begin{corollary}
    A homomorphism of fields is always injective.
\end{corollary}

\begin{proof}
    Let $\varphi : \K \to \L$ be a field homomorphism. Then $\ker \varphi$ is an ideal in $\K$. We have $1_\K \notin \ker \varphi$, therefore $\ker \varphi = \{ 0 \}$ and $\varphi$ is injective.
\end{proof}

\subsection{Field extensions}

We can construct new fields by considering subfields generated by a subset of a field.

\begin{lemma}
    Suppose $\K$ is a field and $\{ \F_\alpha : \alpha \in I \}$ is a family of subfields of $\K$. Then
    \[
        \bigcap_{\alpha \in I} F_\alpha
    \]
    is a subfield of $\K$.
\end{lemma}

\begin{definition}
    Suppose $\L_0 \subseteq \L$ are fields and $\alpha_1, \alpha_2, \ldots \in \L$. Then we define $\L_0(\alpha_1, \alpha_2, \ldots)$ as the intersection of all subfields of $\L$ that contain $\L_0$ and $\alpha_1, \alpha_2, \ldots$, and will also be called the \emph{extension} of $\L_0$ generated by $\alpha_1, \alpha_2, \ldots$.
\end{definition}

It is clear that such an extension is the smallest subfield of $\L$ that contains $\L_0$ and $\alpha_1, \alpha_2, \ldots$.

\begin{example}{\ }
    \begin{enumerate}
        \item Consider the fields $\Q \subseteq \R$, then one can show
        \[
            \Q(\sqrt{2}) = \{ p + q \sqrt{2} : p, q \in \Q \}.
        \]

        \item Since the element $1_\Q$ already generates $\Q$ as a field, it follows that $\Q$ has no proper subfields, i.e. $\Q$ is a \emph{prime field}.
        
        \item The field $\F_p$ is a primary field for any $p \in \P$.
    \end{enumerate}
\end{example}

\begin{definition}
    Let $\F$ be a field. The \emph{primary subfield} of $\F$ is the subfield generated by the element $1_\F$.
\end{definition}

\begin{definition}
    Let $R$ be a ring, then its \emph{characteristic} is defined by
    \[
        \characteristic R \coloneqq \begin{cases}
            \min \{ n \in \N : \sum_{j=1}^n 1_R = 0 \}, & \textrm{if it exists}, \\
            0, & \textrm{otherwise}.
        \end{cases}
    \]
\end{definition}

\begin{lemma}
    Let $\F$ be a field.
    \begin{itemize}
        \item If $\characteristic \F = 0$, then the primary subfield of $\F$ is\footnote{We always mean ``up to isomorphism''.} $\Q$.
        \item If $\characteristic \F = p \in \P$, then the primary subfield of $\F$ is $\F_p$.
    \end{itemize}
\end{lemma}

\begin{lemma}
    Suppose $M \subseteq R$ is a maximal ideal. Then the quotient $R/_M$ is a field.
\end{lemma}

\begin{remark}
    Suppose $\K$ is a field and consider $\K[x]$, the ring of polynomials in $x$ with coefficients in $\K$. Recall that $\K[x]$ is a principle ideal domain. A non-trivial ideal $(f(x))$ is maximal in $\K[x]$ if and only if $f(x)$ is an irreducible polynomial. In particular, $\K[x]/_{(f(x))}$ is a field whenever $f(x)$ is irreducible over $\K$.

    Moreover, $\K[x]/_{(f(x))}$ contains $\K$ as a subfield, via the homomorphisms
    \[
        \K \hookrightarrow \K[x] \to \K[x]/_{(f(x))}.
    \]
    Therefore, for every irreducible $f(x)$, $\K[x]/_{(f(x))}$ is a field extension of $\K$.

    Note that $f(x)$ being irreducible with $\deg f > 1$ implies that $f(x)$ has no roots in $\K$. But $f(x)$ always has a root in $\K[x]/_{(f(x))}$, because $f(x)$ is zero as an element of $\K[x]/_{(f(x))}$.
\end{remark}

The above demonstrates (roughly) that there are relationships:
\begin{align*}
    \textrm{Polynomials} \quad &\leftrightsquigarrow \quad \textrm{Field extensions} \\
    \textrm{Properties of roots} \quad &\leftrightsquigarrow \quad \textrm{Symmetries of field extensions}
\end{align*}

Practically, we want some criteria to determine if a polynomial is irreducible.

\begin{remark}
    If $\K$ is a field and $p(x) \in \K[x]$, $\deg p > 1$ is a polynomial with a root in $\K$, then $p$ is not irreducible. Indeed, for any $a \in \K$ one can write
    \[
        p(x) = q(x) (x-a) + p(a),
    \]
    for some polynomial $q(x) \in \K[x]$. If $a$ is a root of $p(x)$, this demonstrates that $p(x)$ is not irreducible.

    Note that the converse fails: The polynomial $(x^2 + 1)^2$ has no roots in $\R$ and is not irreducible.
\end{remark}

\begin{lemma}
    For a field $\K$, a polynomial $p(x) \in \K[x]$ such that $2 \leq \deg p \leq 3$, is irreducible if and only if it has no roots in $\K$.
\end{lemma}

\begin{lemma}[Gauss]
    Suppose $p(x) = a_0 + a_1 x + \ldots + a_n x^n \in \Z[x]$ is \emph{primitive}, that is $\gcd(a_0, a_1, \ldots, a_n) = 1$. Then $p(x)$ is irreducible in $\Z[x]$ if and only if it is irreducible in $\Q[x]$.
\end{lemma}

\begin{lemma}[Eisenstein's Criterion]
    Let $q(x) = a_0 + a_1 x + \ldots + a_n x^n \in \Z[x]$ and suppose there exists a prime $p \in \P$ such that:
    \begin{enumerate}
        \item $p \nmid a_n$
        \item $p \mid a_0, \ldots, a_{n-1}$
        \item $p^2 \nmid a_0$
    \end{enumerate}
    Then $q(x)$ is irreducible in $\Q[x]$.
\end{lemma}

\begin{example}{\ }
    \begin{enumerate}
        \item The polynomial
        \[
            q(x) = 2x^3 + 6x^2 - 9x + 15
        \]
        satisfies Eisenstein's criterion with $p=3$.

        \item The polynomial
        \[
            q(x) = x^2 + x + 1
        \]
        does not satisfy Eisenstein's criterion. Nevertheless, one can consider the transformation
        \[
            r(x) = q(x+1) = (x+1)^2 + (x+1) + 1 = x^2 + 3x + 3,
        \]
        which does satisfy the criterion. Hence $q(x)$ is also irreducible.
        
        More generally, using the same method as above, one can show that for any prime $p \in \P$, the polynomial
        \[
            q(x) = x^{p-1} + x^{p-2} + \ldots + 1
        \]
        is irreducible.
    \end{enumerate}
\end{example}